%======================================================================
%  CRESCO8 as the experimental testbed
%======================================================================

\section{CRESCO8 as the experimental testbed}
\label{sec:cresco8}

\subsection{Documented configuration}

CRESCO8 is the current CRESCO system at the ENEA research centre in
Portici.  It entered service in 2025 and is connected by a 200\,Gb/s
Mellanox NDR fabric~\cite{cresco8doc,eneapress}.  The operational user
page distinguishes the three compute-node classes listed in
\tabref{tab:cresco8-node-classes}:

\begin{table}[htbp]
\centering
\small
\tabsetup
\caption{CRESCO8 node classes documented by ENEAGRID.  Counts describe a
dated operational document, not a permanent scheduler inventory
~\cite{cresco8doc}.}
\label{tab:cresco8-node-classes}
\begin{tabularx}{\textwidth}{@{}lcccY@{}}
\toprule
\headrow
\textbf{Class} & \textbf{Nodes} & \textbf{CPU cores/node}
 & \textbf{Memory/node} & \textbf{Processor or accelerator} \\
\midrule
CPU
 & 760 & 128 & 512\,GB DDR5
 & 2$\times$ Xeon Platinum 8592+ \\
CPU--HBM
 & 15 & 112 & 1,024\,GB DDR5 + 128\,GB HBM
 & 2$\times$ Xeon CPU Max 9480 \\
GPU
 & 17 & 128 & 512\,GB DDR5 + accelerator HBM
 & 2$\times$ Xeon Platinum 8592+ and
   4$\times$ Intel Data Center GPU Max 1550 \\
\bottomrule
\end{tabularx}
\end{table}

The same page calls this a 793-node system, although the listed classes
sum to 792 and its reported core total is consistent with the 792-node
breakdown.  Earlier commissioning material gives yet another snapshot
~\cite{eneacresco8slides}.  This inconsistency is operational rather than scientific, and the
response is procedural: every experiment retains dated
\texttt{sinfo -N} and \texttt{scontrol show node} output, and the thesis
\textbf{reports the observed inventory} rather than silently reconciling
historical documents.


\subsection{Why the HBM partition supports a clean comparison}

Each documented HBM node contains two Max 9480 sockets, 1\,TB of DDR5
and 128\,GB of HBM.  This makes three comparisons possible, but they do
not have equal evidential strength:

\begin{enumerate}
  \item \textbf{Local HBM versus local DDR5 on the same Max 9480
        socket} is the primary comparison.  Processor, core count,
        executable and most firmware conditions can be held constant.

  \item \textbf{One socket versus two sockets, or Quadrant versus SNC4},
        tests the runtime's response to topology.  It is a placement
        experiment, not a pure memory-technology comparison.

  \item \textbf{Max 9480 versus Platinum 8592+} is a secondary
        system-level baseline.  The processors differ in generation,
        core count, cache, clocks, memory channels and uncore behaviour.
        The contrast cannot identify an HBM effect by itself.
\end{enumerate}

The installed DDR capacity does not reveal the DIMM layout.  Sixteen
64\,GB DIMMs would populate all channels, but the platform also permits
eight 128\,GB DIMMs with Xeon Max~\cite{lp1603}.  Channel population,
memory frequency and local bandwidth therefore require direct
inspection.  The presence of DDR5 also means that the nodes use Flat or
Cache mode rather than the no-DDR HBM-only configuration, but the ENEA
page does not identify which mode is active.  In Flat mode, HBM appears
as memory-only NUMA nodes; in Cache mode it does not appear separately.

The NDR link has a raw one-direction rate of 25\,GB/s per node, far
below local HBM bandwidth.  Multi-node inference may still be necessary
for capacity, but communication can then dominate the placement effect.
Single-node experiments should establish a stable result before adding
distributed parallelism.

\subsection{Scheduler and storage context}

The documented production partition is \texttt{cresco8\_hbm}, with 14
nodes and a 24-hour limit; \texttt{cresco8\_hbm\_dbg} lists two nodes
and a one-hour limit~\cite{cresco8doc}.  Their hostnames can overlap, so
these counts must not be added.

Compute nodes are reached only through Slurm.  The front ends are for
editing, compilation and submission, not benchmarking.  ENEAGRID
provides AFS for shared home directories and Lustre for parallel I/O;
large model files and benchmark output belong on the latter, subject to
the site's current mount points and quotas.  Compiler, MPI and profiler
modules listed on the web page are useful starting points, but an
experiment must record the modules and linked libraries actually loaded
on the run date.

\subsection{Why this platform, and what it costs}
\label{sec:why-cresco8}

Machines that expose HBM and DDR as separately addressable tiers on a
general-purpose CPU are not common, and the ones that exist are mostly
not reachable.  Aurora is the largest Xeon Max deployment and the source
of most of the published measurement literature on this
processor~\cite{ibeid2025aurora,goto2026aurorahpl}, but access to it is
not something a thesis can assume.  Fifteen nodes at ENEA are a small
resource by comparison, and for the question in
\secref{sec:research-gap} that hardly matters: the primary
comparison is defined on a single socket, and a single socket is what
the argument needs.  The scarce commodity here is \emph{not nodes but the processor
itself}.

The scarcity does shape the experimental design, though, and it is
better to admit that in the method than to discover it halfway through.
Fourteen nodes in \texttt{cresco8\_hbm} are shared with the rest of the
centre, the queue limit is 24 hours, and the debug partition offers two
nodes for an hour.  A full Cartesian sweep over model, quantization,
prompt length, generation length, batch size, thread count, memory tier
and clustering mode does not fit in that budget and would not be worth
running if it did.  Salaria et al.\ give the principled response:
measure the cost of the benchmark alongside its result, and choose a
subset of the space that preserves most of its information.  Their
example reduces generation length from 1024 to 128 tokens and retains
96.6\% of the accuracy at 2.7$\times$ less
cost~\cite{salaria2025metametrics}.  \secref{sec:inference-model}
supplies the other half of the argument, since the model predicts which
factors move the outcome and which merely enlarge the table.

One further consequence deserves stating early.  Because HBM nodes are
shared and the memory mode is a firmware setting fixed at boot, the mode
available on the allocated nodes is a constraint rather than a variable.
If the nodes are found in Cache mode and the administrators cannot
provide Flat mode, the primary comparison in this thesis is not
possible as specified, and the design has to change rather than the
claim.  Establishing the active mode is therefore the first task of the
first allocation, and not a detail to be confirmed later.

\subsection{First allocation: capture facts before benchmarking}
\label{sec:first-allocation}

The first HBM allocation should create a machine-readable system
manifest.  The following script is deliberately a topology-capture
skeleton rather than a performance benchmark; account names, available
utilities and Slurm resource syntax must be checked on the live system.

\begin{lstlisting}[style=shell,
  caption={Minimal exclusive allocation for a CRESCO8 topology
  manifest.}]
#!/bin/bash
#SBATCH --nodes=1
#SBATCH --ntasks=1
#SBATCH --cpus-per-task=112
#SBATCH --partition=cresco8_hbm
#SBATCH --exclusive
#SBATCH --hint=nomultithread
#SBATCH --output=%j-system.txt

set -u

echo "job_id=${SLURM_JOB_ID}"
echo "host=$(hostname)"
date --iso-8601=seconds
scontrol show job "${SLURM_JOB_ID}"
lscpu
lscpu --extended
numactl --hardware
lstopo-no-graphics 2>/dev/null || true
cat /sys/devices/system/cpu/smt/active 2>/dev/null || true
findmnt -t lustre 2>/dev/null || true
module list 2>&1 || true
\end{lstlisting}

\figref{fig:cresco8-lstopo} is what the topology half of that manifest
actually looks like, captured from a CRESCO8 node.  It is worth reading
once, because it shows the difference between a specification and an
observation.  The two processor packages appear with their own memory,
their caches and the devices attached to each, and the numbers on the
links are relative access costs.  Note that this particular node is from
the general CPU partition rather than the HBM partition: it reports two
memory nodes of roughly 250\,GB each and no separate high-bandwidth
tier.  On a Max~9480 node running in Flat mode, the same command should
additionally report memory-only NUMA nodes carrying the HBM, and
confirming that is the first thing \secref{sec:hbm-modes} requires of
the first allocation.

\begin{figure}[htbp]
\centering
\includegraphics[width=\textwidth]{cresco8-cpu-node-diagram.png}
\caption{Topology of a CRESCO8 compute node as reported by
\texttt{lstopo}, captured on \texttt{cresco8x001.portici.enea.it}.  Two
packages, 64 cores each, with their own memory and network devices.  The
absence of memory-only NUMA nodes identifies this as a node of the
general CPU partition; an HBM node in Flat mode would show them.}
\label{fig:cresco8-lstopo}
\end{figure}

The manifest should be extended, subject to permissions, with the
following evidence:

\begin{table}[htbp]
\centering
\small
\tabsetup
\caption{Minimum platform record for an inference result.}
\label{tab:platform-record}
\begin{tabularx}{\textwidth}{@{}lY@{}}
\toprule
\headrow
\textbf{Category} & \textbf{Record} \\
\midrule
Identity
 & Hostname, Slurm job ID, partition, date and node allocation. \\
Processor
 & Exact SKU, sockets, physical cores, SMT state, microcode, governor,
   turbo policy and observed frequency. \\
Memory
 & Memory and clustering modes, NUMA distance matrix, HBM and DDR
   node IDs, DIMM population, huge-page state and measured local/remote
   bandwidth. \\
I/O
 & Lustre mount and striping used for model loading; NIC identity and
   NUMA affinity if multi-node tests are attempted. \\
Software
 & OS and kernel, loaded modules, compiler, runtime commit, linked
   libraries, build flags and complete command line. \\
Telemetry
 & Available RAPL, performance-counter and management-controller
   interfaces, including permissions and measurement boundaries. \\
\bottomrule
\end{tabularx}
\end{table}

This separation between documented configuration and observed
configuration is important.  Vendor documents establish what the platform \emph{can} support, and
ENEAGRID describes a dated cluster state.  \textbf{Only the archived
manifest establishes the conditions under which a reported result was
actually obtained.}
